\section{A model of arbitrage under India's limits}\label{sec:model}

Before looking at any data I wrote a model in which each Indian limit to arbitrage enters as a parameter. Its purpose is
to say in advance what each limit should do, and what result would contradict it. This section presents the model and
the propositions it implies; \cref{app:proofs} gives the proofs.

\subsection{Set-up}
Let $S$ be the price of a stock and $F$ the price of its future, which expires at $T$. With an interest rate $r$ and
dividends whose present value is $\mathrm{PV}(D)$, the cost-of-carry value of the future is
$F^*=(S-\mathrm{PV}(D))\,e^{r(T-t)}$, and the basis is
\begin{equation}\label{eq:basis}
m=\ln F-\ln F^* ,
\end{equation}
measured in basis points of the stock price. A positive basis means the future is rich relative to the stock.

Two groups trade. \emph{Outside traders}, such as speculators and hedgers, trade for their own reasons. Their net demand
in the cash market is $D_S-\psi_S s$ and in the futures market $D_F-\psi_F f$, where $s$ and $f$ are the deviations of
the stock and futures prices from fundamental value (so $m\approx f-s$), $D_S$ and $D_F$ shift demand, and $\psi_S$ and
$\psi_F$ measure how deep each market is. An \emph{arbitrage sector}, standing for arbitrage funds and proprietary
desks, holds $a$ units of the cash-and-carry trade: long the stock and short the future ($a<0$ would be the reverse
trade). Market clearing in both markets gives
\begin{equation}\label{eq:clear}
m=\frac{u-a}{\Phi},\qquad \frac{1}{\Phi}=\frac{1}{\psi_F}+\frac{1}{\psi_S},\qquad
u=\Phi\Big(\frac{D_F}{\psi_F}-\frac{D_S}{\psi_S}\Big).
\end{equation}
The variable $u$ is the \emph{relative demand pressure} of outside traders toward the future: it is high when they want
more exposure through futures than through the stock, for example when speculators seek leverage. Without arbitrage
the basis would be $u/\Phi$; every unit of cash-and-carry lowers it by $1/\Phi$.

The arbitrage sector chooses $a$ to maximise
\begin{equation}\label{eq:objective}
a\,m-\frac{\gamma}{2}a^2-c_+a^+-c_-a^-,\qquad \gamma=\alpha\sigma^2 ,
\end{equation}
where $a^\pm=\max(\pm a,0)$. The first term is what the position earns when the basis closes at expiry. The second is
the cost of risk: $\alpha$ is risk aversion per unit of capital, which is higher when arbitrage capital is scarce, and
$\sigma^2$ is the risk that the basis widens before it closes. The last two terms are the costs of each trade: $c_+$
for cash-and-carry (taxes, fees, price impact, and any gap between the return arbitrage capital requires and the rate
used to compute the basis) and $c_-$ for the reverse trade, which also pays for borrowing the stock and which is very
high in India because overnight short sales are hard. \Cref{tab:params} lists how each Indian limit enters.

\begin{table}[t]
\caption{Each Indian limit to arbitrage enters the model as a parameter}\label{tab:params}
\exsub{Limits, their Indian form, and where they appear in the model}
\small\begin{tabularx}{\linewidth}{@{}K{3.0cm}L>{\raggedright\arraybackslash}p{3.3cm}@{}}
\toprule
\thd{Limit} & \thd{Indian form} & \thd{Parameter} \\
\midrule
Trading costs & STT on futures and on delivery, stamp duty, fees, brokerage & Costs $c_+$, $c_-$ \\ \rowrule
Short-sale constraints & Shorting only in F\&O stocks; thin stock lending; the overnight short is a future & Cost of reverse trade $c_-$ \\ \rowrule
Arbitrage capital & Arbitrage-fund assets; funding rates; margin rules & Risk aversion $\alpha$, so $\gamma$ \\ \rowrule
Arbitrage risk & Volatility of the basis before it closes & Risk $\sigma^2$, so $\gamma$ \\ \rowrule
Position limits & MWPL; the F\&O ban (enter above 95\%, leave at 80\%) & Reduce-only rule (\cref{sec:model-ban}) \\ \rowrule
Market depth & Liquidity of outside traders in each market & $\psi_S$, $\psi_F$, so $\Phi$ \\
\bottomrule
\end{tabularx}
\exsource{Source: the model of this section.}
\end{table}

\subsection{The no-arbitrage band}
\begin{proposition}[An asymmetric band]\label{prop:band}
The equilibrium basis is unique and equals
\begin{equation}\label{eq:band}
m(u)=\begin{cases}
\dfrac{\gamma u + c_+}{\gamma\Phi+1}, & u>\Phi c_+,\\[2mm]
\dfrac{u}{\Phi}, & -\Phi c_-\le u\le \Phi c_+,\\[2mm]
\dfrac{\gamma u - c_-}{\gamma\Phi+1}, & u<-\Phi c_- .
\end{cases}
\end{equation}
\end{proposition}
In words: as long as the basis lies between $-c_-$ and $c_+$, no arbitrage pays, and the basis simply follows the demand
of outside traders. Once it moves beyond an edge, arbitrage starts and absorbs part of each further push, so the basis
moves less (\cref{fig:modelband}). Because the reverse trade is costlier, the band reaches further below zero than above
it.

\begin{figure}[t]
\caption{Inside the band the basis follows demand one for one; outside it arbitrage absorbs part of each push}\label{fig:modelband}
\exsub{Equilibrium basis against the demand pressure of outside traders, with illustrative costs $c_+=15$\bp\ and $c_-=60$\bp}
\centering\includegraphics[width=0.82\linewidth]{model_band.pdf}
\exsource{Source: \cref{prop:band}; illustrative parameters, not estimates.}
\end{figure}

\begin{proposition}[What moves the band]\label{prop:cs}
(i) The edges of the band are $c_+$ and $-c_-$: costs set the edges, and a change in a tax moves the corresponding edge
one for one. (ii) Capital and risk do not move the edges; they set how far the basis travels beyond them, since the
slope $\gamma/(\gamma\Phi+1)$ rises with $\gamma$. (iii) The carrying part of the costs is proportional to the time left
to expiry, so the band narrows as expiry approaches. (iv) The band is centred at $(c_+-c_-)/2$, below zero whenever the
reverse trade costs more than cash-and-carry.
\end{proposition}
Part (i) gives a direct test. When the STT on futures was raised, the cost of cash-and-carry rose, so the upper edge of
the band should have risen, and by more for larger increases.

\subsection{Existing positions}
Arbitrageurs do not start each day from zero. Suppose the sector enters the day holding $a_{t-1}$ units and pays $c_u$ per
unit to unwind them.

\begin{proposition}[The band moves with the existing book]\label{prop:inv}
The sector adds to its position only if $m_t>\gamma a_{t-1}+c_+$, unwinds only if $m_t<\gamma a_{t-1}-c_u$, and otherwise
holds its position. The band within which it does not trade is $[\gamma a_{t-1}-c_u,\ \gamma a_{t-1}+c_+]$.
\end{proposition}
The edges of the band are therefore not fixed: they shift with the size of the existing book. A threshold model with fixed
thresholds, the standard tool, will estimate an average position of the edges rather than the costs themselves. I
checked this before using real data: when I simulated the basis from this model with costs of 15\bp\ to add and 10\bp\
to unwind, the threshold model placed the thresholds at about $-10$ and $0$\bp, so the estimated upper threshold lay well
below the 15\bp\ cost of adding to a position. This is why my test of costs compares how
the edges \emph{change} around tax increases rather than comparing their level with the costs.

\subsection{The F\&O ban}\label{sec:model-ban}
The ban acts on everyone in the market, so the model represents it with two restrictions.

\begin{assumption}[The ban]\label{as:ban}
During a ban, (B1) the arbitrage sector may only reduce its position, $0\le a_t\le a_{t-1}$; and (B2) outside traders
cannot open new futures positions. A share $\chi$ of this blocked demand moves to the stock (a trader who cannot buy the
future buys the share instead) and the rest is withdrawn. Outside traders can still close existing positions.
\end{assumption}

\begin{proposition}[One-sided correction]\label{prop:onesided}
During a ban, rich futures are no longer corrected: since nobody may open new cash-and-carry positions, any further
pressure toward the future passes through to the basis in full. Cheap futures are still corrected, by arbitrageurs
unwinding their positions, as long as they have positions left to unwind. This correction weakens as positions run
down, and an expiry inside the ban removes the expiring positions, since they cannot be rolled.
\end{proposition}

\begin{proposition}[The level effect has no fixed sign]\label{prop:level}
Compared with normal times, a ban raises the basis when the pressure on a stock comes mainly from traders closing
positions, and lowers it when new demand keeps arriving, because part of the blocked demand for futures moves to the
stock. Which case applies depends on the composition of trading before the ban.
\end{proposition}

\begin{proposition}[Trading in a ban]\label{prop:trade}
If the rule binds on every trade, each futures trade during a ban closes a long position and a short position. The
change in open interest then equals minus the volume traded, so the ratio $R=(\Delta OI+V)/V$, which is close to one on
normal days, is zero, and futures volume falls. Trades that offset positions across instruments, which are allowed from
October 2025, break this identity.
\end{proposition}

\Cref{fig:modelpath} shows a simulated ban. When the ban starts, the arbitrage book can only shrink, and the basis follows
demand upward without resistance. One further implication concerns the days before a ban. If closing a position is
costlier during a ban, an arbitrageur who opens a position when utilisation is already high must allow for the chance
of being caught in one, so the upper edge of the band should rise as utilisation approaches 95\%. The model does not,
however, say whether utilisation should bunch just below the trigger: arbitrageurs holding back would make it thinner
there, while speculators rushing in before a ban would do the opposite.

\begin{figure}[t]
\caption{In a ban the arbitrage book can only shrink, so upward pressure on the basis passes through in full}\label{fig:modelpath}
\exsub{One simulated path with and without a ban on days 120 to 159, same shocks; basis (top) and arbitrage position (bottom)}
\centering\includegraphics[width=0.86\linewidth]{model_path.pdf}
\exsource{Source: the model of this section, simulated; illustrative parameters, not estimates.}
\end{figure}

\subsection{From the model to the tests}
\Cref{tab:predictions} lists the predictions, the result that would contradict each, and where it is tested. The
predictions about the band, about costs and frictions, and about trading in a ban are tested in this report. The
predictions about what the ban does to the level and the correction of the basis, and the model's predictions about
capital and short-sale constraints, are tested in the second stage. The labels in the last column are those of the
pre-analysis plan, which names four primary hypotheses (H3, H6, H7, H8) and corrects them jointly for multiple testing.

\begin{table}[t]
\caption{Every prediction comes with the result that would contradict it}\label{tab:predictions}
\exsub{Predictions of the model, and where each is tested; primary hypotheses of the test plan in bold}
\small\begin{tabularx}{\linewidth}{@{}P{5.2cm}P{3.4cm}P{2.6cm}L@{}}
\toprule
\thd{Prediction} & \thd{Contradicted if} & \thd{Tested in} & \thd{Plan label} \\
\midrule
\multicolumn{4}{@{}l}{\thd{This report}} \\
The basis moves in an asymmetric band, wider below zero (Props.\ 1, 2) & No band; a symmetric band & \Cref{sec:band} & Band estimation \\ \rowrule
Tax increases raise the upper edge, more for larger increases (Prop.\ 2) & Edges unrelated to cost changes & \Cref{sec:costs} & \textbf{H3} \\ \rowrule
Bands are wider where frictions are higher (Props.\ 1 to 3) & No relation & \Cref{sec:costs} & H1 \\ \rowrule
In a ban futures volume falls and $\Delta OI\approx -V$ (Prop.\ 6) & Futures activity unchanged & \Cref{sec:ban} & \textbf{H8} \\ \rowrule
Utilisation near the 95\% trigger (not signed) & (descriptive) & \Cref{sec:ban} & Density test \\
\midrule
\multicolumn{4}{@{}l}{\thd{Second stage}} \\
In a ban rich futures revert more slowly (Prop.\ 4) & No asymmetry, or the opposite & Stage II & \textbf{H7} \\ \rowrule
The level effect follows the composition of trading (Prop.\ 5) & Sign unrelated to composition & Stage II & Level effect \\ \rowrule
The upper tail rises as utilisation nears 95\% & A jump only at 95\% & Stage II & H9 \\ \rowrule
Scarcer capital lets the basis travel further (Prop.\ 2) & No relation & Stage II & H4 \\ \rowrule
Access to short selling speeds price adjustment and lowers overpricing & No change, or the opposite & Stage II & \textbf{H6}, H5, H2 \\
\bottomrule
\end{tabularx}
\exsource{Source: the model of this section; the pre-analysis plan (\cref{app:plan}).}
\end{table}
